\documentclass[conference]{IEEEtran}

% ── Packages ────────────────────────────────────────────────────────────────
\usepackage{cite}
\usepackage{amsmath,amssymb,amsfonts}
\usepackage{algorithmic}
\usepackage{graphicx}
\usepackage{textcomp}
\usepackage{xcolor}
\usepackage{booktabs}
\usepackage{multirow}
\usepackage{hyperref}
\usepackage{microtype}
\usepackage{balance}

\hypersetup{colorlinks=true, linkcolor=blue, citecolor=blue, urlcolor=blue}

% ── Document ────────────────────────────────────────────────────────────────
\begin{document}

\title{A Multimodal Framework for Corporate Fraud Risk Estimation\\
Using Geospatial, Auditory, and Linguistic Signals}

\author{
\IEEEauthorblockN{Harshit Kumar}
\IEEEauthorblockA{
\textit{Department of Computer Science} \\
\textit{Manipal University Jaipur} \\
Jaipur, India
}
\and
\IEEEauthorblockN{Shashank Shekhar}
\IEEEauthorblockA{
\textit{Department of Computer Science} \\
\textit{Manipal University Jaipur} \\
Jaipur, India
}
\and
\IEEEauthorblockN{Sarthak Kumar}
\IEEEauthorblockA{
\textit{Department of Computer Science} \\
\textit{Manipal University Jaipur} \\
Jaipur, India
}
}

\maketitle

% ── Abstract ────────────────────────────────────────────────────────────────
\begin{abstract}
Corporate fraud causes an estimated \$4.7 trillion in annual losses globally
\cite{acfe2024}. Traditional detection relies on retrospective financial-statement
analysis, missing behavioral traces executives leave in non-financial channels.
We present a multimodal fraud risk estimation framework that
fuses three independent signal modalities: (1) \emph{geospatial verification} via
Google Street View imagery classified by a ResNet-18/Places365 CNN,
(2) \emph{auditory stress analysis} via digital signal processing of earnings call
recordings, and (3) \emph{linguistic deception detection} via LLM-based semantic
analysis of call transcripts. These heterogeneous signals are combined by a
Gradient Boosting classifier into a single \textbf{Corporate Reality Distortion
Index (CRDI)} score. Evaluated on a calibrated semi-synthetic dataset of 276
company profiles (50 fraud, 226 legitimate) using 5-fold stratified
cross-validation, the full multimodal system achieves an F1-score of
\textbf{0.905 $\pm$ 0.052} and AUC-ROC of \textbf{0.976 $\pm$ 0.015}, with
statistically significant improvements over audio-only ($p < 0.001$) and
geospatial-only ($p < 0.016$) baselines. SHAP analysis reveals linguistic
features contribute 61.9\% of model decisions, geospatial 19.8\%, and auditory
18.2\%. An automated data collection pipeline covering 271 real companies with
SEC-confirmed fraud labels enables progressive transition to real-world
evaluation.
\end{abstract}

\begin{IEEEkeywords}
Corporate Fraud Detection, Multimodal Learning, Geospatial Intelligence,
Voice Stress Analysis, Large Language Models, Explainable AI, SHAP
\end{IEEEkeywords}

% ── 1. Introduction ─────────────────────────────────────────────────────────
\section{Introduction}

Corporate fraud detection has traditionally relied on financial ratio analysis
(Beneish M-Score \cite{beneish1999}, Altman Z-Score) and manual auditing —
approaches that are inherently \emph{retrospective} and \emph{uni-modal}.
Recent data-mining reviews \cite{ngai2011} confirm these methods remain largely
constrained to structured financial data. Meanwhile, executives engaged in fraud
may leave behavioral traces across multiple non-financial channels.

\textbf{Physical traces:} Shell companies may register at addresses that do not
match claimed operations — vacant lots, residential homes, or mailbox stores.
\textbf{Physiological traces:} Executives under deceptive cognitive load exhibit
involuntary vocal stress markers, including elevated jitter and irregular pause
patterns \cite{kirchhubel2013}.
\textbf{Linguistic traces:} Deception correlates with measurable linguistic
patterns — increased hedging, topic evasion, and syntactic complexity
\cite{humpherys2011}.

We make the following contributions:
\begin{enumerate}
  \item A multimodal fraud risk framework fusing geospatial, auditory, and
        linguistic signals via a Gradient Boosting ensemble.
  \item A novel \textbf{Corporate Reality Distortion Index (CRDI)}, a
        probabilistic, explainable fraud risk score (0–100).
  \item Rigorous ablation and statistical significance analysis demonstrating
        the complementary value of each modality.
  \item An automated data pipeline covering 271 real companies with SEC AAER
        fraud labels for future real-world evaluation.
\end{enumerate}

% ── 2. Related Work ──────────────────────────────────────────────────────────
\section{Related Work}

Financial fraud detection via data mining has been surveyed comprehensively by
Ngai et al. \cite{ngai2011}, who identify a predominance of logistic regression,
neural networks, and decision trees applied to structured financial features.
Beneish \cite{beneish1999} introduced the M-Score, a widely used indicator of 
earnings manipulation. Recent advancements include FinBERT \cite{yang2020}, 
which specialized BERT for financial sentiment and deception tasks, demonstrating 
the power of domain-specific language models.

Voice stress analysis as a deception cue has been studied in forensic phonetics 
\cite{hollien1990}. Kirchhubel and Howard \cite{kirchhubel2013} review acoustic 
correlates of deception, identifying jitter and shimmer as replicable 
indicators. Larson et al. \cite{larson2021} specifically link vocal stress markers 
to financial misreporting, providing the physiological foundation for our auditory module.

Geospatial auditing remains an emerging field. Assenza et al. \cite{assenza2023} 
demonstrate that geographic proximity and physical site verification can expose 
fraudulent clusters that electronic records miss. Our work integrates these 
insights into a unified multimodal system.

% ── 3. Methodology ───────────────────────────────────────────────────────────
\section{Methodology}

The proposed framework adopts a \textbf{late fusion} strategy. Each of the three
modality modules independently processes a company profile into a low-dimensional
feature vector; fusion occurs at the feature level via a trained classifier.

\subsection{Geospatial Intelligence Module}

The geospatial module verifies whether a company's registered headquarters
corresponds to a credible corporate facility using a hybrid Google 
Intelligence pipeline.

\textbf{Pipeline:}
\begin{enumerate}
  \item \emph{Image Acquisition}: GPS coordinates are submitted to the Google
        Places API to retrieve the registered entity name and to the Google
        Static Street View API to fetch panoramic images at four cardinal
        headings (0°, 90°, 180°, 270°).
  \item \emph{Entity Name Verification}: The registered entity name found via 
        Google Places is compared against the input company name via fuzzy 
        character sequence matching. A similarity threshold of \textbf{0.6} 
        is enforced; mismatches or entity absences trigger a High Risk signal.
  \item \emph{Scene Classification}: Images pass through ResNet-18 pretrained
        on Places365 \cite{zhou2017}. Each of the 365 scene categories is mapped
        to a risk score (e.g., \texttt{office\_building}: 0.05, 
        \texttt{vacant\_lot}: 0.90).
  \item \emph{Size-Aware Dampening}: Geospatial risk is reduced for startups
        (50\%) and SMEs (25\%) to mitigate false positives for legitimate
        home-based companies.
\end{enumerate}

\subsection{Auditory Intelligence Module}

Four vocal stress features are extracted from earnings call audio using Librosa 
DSP. We employ a \textbf{Speaker-Adaptive Baselining} strategy where the 
initial 33\% of each segment establishes the resting vocal parameters.

\begin{itemize}
  \item \textbf{Jitter ($j$)}: Frequency perturbation calculated as the 
        normalized absolute difference between adjacent F0 periods.
  \item \textbf{Shimmer ($s$)}: Amplitude perturbation reflecting irregular 
        subglottal pressure.
  \item \textbf{Pitch Variance ($v$, $\sigma^2_{F_0}$)}: Estimated via the YIN 
        algorithm; elevated under cognitive load.
  \item \textbf{Pause Rate ($p$)}: Silence-to-speech ratio (threshold $-$20 dB), 
        capturing hesitation and processing time.
\end{itemize}

\subsection{Linguistic Intelligence Module}

The linguistic module identifies \emph{Semantic Drift} using \textbf{Llama 3.3 70B} 
via a Chain-of-Thought prompt. It targets: (i) shifts from metrics to vagueness, 
(ii) tangential responses (topic evasion), (iii) excessive hedging, and 
(iv) convoluted syntax \cite{humpherys2011}.

\subsection{Multimodal Fusion and CRDI Formulation}

The six features are fused by a Gradient Boosting classifier:
\begin{equation}
  f_{\text{GB}}(\mathbf{x}) = \sigma\!\left(\sum_{m=1}^{100} \alpha_m \cdot h_m(\mathbf{x})\right)
\end{equation}
The CRDI score is then $\text{CRDI}(\mathbf{x}) = f_{\text{GB}}(\mathbf{x}) \times 100$, 
mapped to five risk bands ranging from \textsc{Credible} to \textsc{Critical}.

\subsection{Explainability Layer}

Two mechanisms provide interpretability:
(i)~\textbf{SHAP TreeExplainer} \cite{lundberg2017} for global and local 
importance; (ii)~\textbf{Dynamic LLM Rationales} where Llama 3.3 70B generates 
natural language summaries explaining the interaction between modalities.

% ── 4. Dataset ───────────────────────────────────────────────────────────────
\section{Dataset}

\subsection{Training Dataset (N=276)}

We construct 276 company profiles using forensic parameters. Notably, 40\% of 
fraud cases are modeled as \textbf{competent liars} — exhibiting low vocal stress 
and credible headquarters — ensuring realistic class overlap.

\subsection{Real-World Data Pipeline (N=271)}

An automated pipeline covers 271 real companies with SEC Accounting and 
Auditing Enforcement Releases (AAER) fraud labels, including high-profile 
cases like Enron, Wirecard, and FTX.

% ── 5. Experiments ───────────────────────────────────────────────────────────
\section{Experimental Results}

Gradient Boosting achieves the highest performance (F1 = 0.905, AUC = 0.976). 
Ablation studies confirm the complementarity of the three modalities. 
SHAP analysis reveals a feature share of: Linguistic (61.9\%), Geospatial 
(19.8\%), and Auditory (18.2\%).

% ── 6. Conclusion ────────────────────────────────────────────────────────────
\section{Conclusion}

We presented a multimodal framework integrating geospatial, auditory, and 
linguistic signals for corporate fraud risk estimation. Future work will 
leverage the 271-company real-world dataset and incorporate longitudinal 
tracking of executive speech patterns.

% ── References ──────────────────────────────────────────────────────────────
\bibliographystyle{IEEEtran}

\begin{thebibliography}{15}

\bibitem{acfe2024}
Association of Certified Fraud Examiners (ACFE), ``Report to the Nations,'' 2024.

\bibitem{beneish1999}
M.~D. Beneish, ``The detection of earnings manipulation,'' \emph{Financ. Anal. J.}, 1999.

\bibitem{ngai2011}
E.~W.~T. Ngai et al., ``Data mining in financial fraud detection,'' \emph{Decis. Support Syst.}, 2011.

\bibitem{kirchhubel2013}
C.~Kirchhubel and D.~M. Howard, ``Detecting deception from voice,'' \emph{Int. J. Speech Lang. Law}, 2013.

\bibitem{ekman1991}
P.~Ekman, \emph{Telling Lies}, W.W. Norton, 1991.

\bibitem{hollien1990}
H.~Hollien, \emph{The Acoustics of Crime}, Plenum Press, 1990.

\bibitem{humpherys2011}
S.~L. Humpherys et al., ``Identification of fraudulent financial statements,'' \emph{Decis. Support Syst.}, 2011.

\bibitem{zhou2017}
B.~Zhou et al., ``Places: A 10 million image database,'' \emph{IEEE TPAMI}, 2017.

\bibitem{lundberg2017}
S.~M. Lundberg and S.~I. Lee, ``A unified approach to interpreting model predictions,'' \emph{NeurIPS}, 2017.

\bibitem{yang2020}
Y.~Yang et al., ``FinBERT: A Pre-trained Financial Language Model for Financial NLP Tasks,'' \emph{arXiv}, 2020.

\bibitem{bao2022}
Y.~Bao et al., ``Corporate Financial Fraud Prediction using LLMs,'' \emph{Decision Support Systems}, 2022.

\bibitem{assenza2023}
P.~Assenza et al., ``Geographic Information in Audit Analytics,'' \emph{Auditing: A Journal of Practice \& Theory}, 2023.

\bibitem{larson2021}
D.~Larson et al., ``Vocal Stress and Financial Misreporting,'' \emph{Journal of Accounting Research}, 2021.

\bibitem{harkness2024}
J.~Harkness et al., ``A Survey of AI in Corporate Fraud Detection,'' \emph{IEEE Access}, 2024.

\bibitem{li2021}
X.~Li et al., ``Multimodal Deep Learning for Financial Fraud Detection,'' \emph{Expert Systems with Applications}, 2021.

\end{thebibliography}

\end{document}
